\section{System Architecture}
\label{sec:architecture}

\begin{figure}[t]
\centering
\begin{tikzpicture}[
  component/.style={draw, rounded corners=1pt, align=center,
    minimum height=8mm, text width=30mm, inner sep=2pt, font=\footnotesize},
  transfer/.style={draw, rounded corners=1pt, align=center,
    text width=74mm, inner sep=4pt, font=\footnotesize},
  flow/.style={-{Stealth[length=4pt]}, semithick},
  both/.style={{Stealth[length=4pt]}-{Stealth[length=4pt]}, semithick}]
  \node[component] (author) at (-20mm,0) {Desktop Creator};
  \node[component] (sqlite) at (20mm,0) {Local SQLite\\editable itinerary};
  \node[transfer] (file) at (0,-14mm) {\texttt{.tt} file\\Trip ID, revision, schedule, and embedded assets};
  \node[transfer,below=5mm of file] (viewer) {Mobile Viewer\\Decode, validate, and compare revisions};
  \node[component] (records) at ([xshift=-20mm,yshift=-9mm]viewer.south) {Local trip store\\and asset files};
  \node[component] (check) at ([xshift=20mm,yshift=-9mm]viewer.south) {Local traveler\\checklist};
  \draw[both] (author) -- (sqlite);
  \draw[flow] (author.south) -- (author.south |- file.north);
  \draw[flow] (file) -- (viewer);
  \draw[flow] (viewer.south -| records.north) -- (records.north);
  \draw[both] (viewer.south -| check.north) -- (check.north);
\end{tikzpicture}

\caption{The Creator publishes a snapshot. The Viewer stores the imported itinerary and traveler checklist in separate local records.}
\label{fig:delivery}
\end{figure}

\subsection{Shared core and storage}
Both apps use the Dart core for the trip model, file codec, and import and export operations. The core defines ports for storage and the codec \cite{hex}. Desktop screens also call SQLite directly for routine editing. On Android, the shared importer saves trips and checklist records as JSON and keeps binary assets in a file store (Fig.~\ref{fig:delivery}).

The stable design calls for transactional Drift/SQLite storage in the Viewer and atomic asset writes. The importer would validate and stage assets before publishing the schedule. Its transaction would preserve the checklist and carry forward the completion state. Once the traveler receives a file, the exchange can proceed without a server or analytics service. Download, messaging, routing, and map services have their own network requirements.

\subsection{Attachment indexing}
The exporter hashes document and event-photo bytes with SHA-256 \cite{sha}. Records that refer to identical bytes can point to one blob. Each asset descriptor gives its digest, size, and offset in the attachment region. For versioned exports, the importer checks attachment bounds and verifies the hashes of all indexed assets. The mobile store names each asset file by its digest. A re-export can index legacy inline data. The banner has its own field and is outside the size experiment in Section~\ref{sec:evaluation}.
